\begin{ouverture}[geometrie]{13}{Angles et droites parallèles}{Partie 3 · Espace et géométrie}
{Vocabulaire · Bissectrice · Angles alternes internes et correspondants}
\objectif{Reconnaître des angles dans une configuration.}
\objectif{Calculer avec des angles adjacents et opposés.}
\objectif{Utiliser les angles pour caractériser un parallélisme.}
\prerequis{Rapporteur et angles usuels.}
\end{ouverture}
\pagestyle{otMain}\essentiel\fontsize{11}{15}\selectfont
\phantomsection\label{cours-13-1}
\essSection{1}{Reconnaître les angles}
\begin{definition}Un angle nul mesure $0^\circ$, droit $90^\circ$, plat $180^\circ$, plein $360^\circ$. Un angle aigu est strictement compris entre 0 et $90^\circ$ ; un angle obtus entre 90 et $180^\circ$. Dans $\widehat{ABC}$, B est le sommet.\end{definition}
\label{fig-cours-13-1-cours-angles}\figureGeo{cours-angles}

\begin{definition}Deux angles adjacents ont le même sommet, un côté commun et des intérieurs disjoints. Deux angles supplémentaires ont une somme de $180^\circ$. La bissectrice partage un angle en deux angles de même mesure.\end{definition}
\begin{propriete}Lorsque deux droites se coupent, deux angles adjacents formant un angle plat ont une somme de $180^\circ$. Deux angles opposés par le sommet ont donc la même mesure : chacun d’eux est supplémentaire d’un même troisième angle, adjacent aux deux.\end{propriete}
\clearpage
\phantomsection\label{cours-13-2}
\essSection{2}{Caractériser les parallèles}
\label{fig-cours-13-2-cours-paralleles}\figureGeo{cours-paralleles}

\begin{definition}Une sécante coupe deux droites. Deux angles alternes internes sont entre les droites et de part et d’autre de la sécante. Deux angles correspondants occupent la même position autour des deux intersections.\end{definition}
\begin{propriete}Si les droites sont parallèles, les angles alternes internes sont égaux, de même que les angles correspondants ; cela se démontre avec le demi-tour (chapitre 15). Réciproquement, l’égalité d’une paire d’angles alternes internes, ou d’une paire d’angles correspondants, prouve le parallélisme ; cette réciproque est admise.\end{propriete}
\begin{exemple}Cas des angles droits : deux droites perpendiculaires à une même droite forment avec elle des angles correspondants droits, donc égaux ; elles sont parallèles. De même, si deux droites sont parallèles, toute perpendiculaire à l’une est perpendiculaire à l’autre.\end{exemple}
\begin{savoirfaire}[Aller des données à la conclusion]Identifier d’abord la paire d’angles. Écrire l’hypothèse connue : droites parallèles ou égalité des angles. Choisir la propriété dans le bon sens puis conclure. Une ressemblance sur le dessin ne suffit pas.\end{savoirfaire}

\clearpage\rubrique{Découverte}
\begin{blocLivre}{ex-13-1}\begin{exercice}[1]{Nommer un angle}
Dans la figure, quel est le sommet de $\widehat{AOB}$ ? L’angle de $40^\circ$ est-il aigu ?
\label{fig-ex-13-1}\figureGeo{ang-1}
\end{exercice}\end{blocLivre}
\begin{blocLivre}{ex-13-2}\begin{exercice}[1]{Un angle droit partagé}
Un angle droit est partagé en deux angles de $35^\circ$ et $x$. Calculer $x$.
\label{fig-ex-13-2}\figureGeo{ang-2}
\end{exercice}\end{blocLivre}
\begin{blocLivre}{ex-13-3}\begin{exercice}[1]{Un angle plat partagé}
Deux angles adjacents forment un angle plat. L’un mesure $128^\circ$. Calculer l’autre.
\label{fig-ex-13-3}\figureGeo{ang-3}
\end{exercice}\end{blocLivre}
\begin{blocLivre}{ex-13-4}\begin{exercice}[1]{Opposés par le sommet}
Deux droites se coupent. Un angle vaut $63^\circ$. Que vaut l’angle opposé par le sommet ?
\label{fig-ex-13-4}\figureGeo{ang-4}
\end{exercice}\end{blocLivre}
\begin{blocLivre}{ex-13-5}\begin{exercice}[1]{Bissectrice}
La demi-droite verte est la bissectrice d’un angle de $74^\circ$. Quelle mesure a chaque moitié ?
\label{fig-ex-13-5}\figureGeo{ang-5}
\end{exercice}\end{blocLivre}

\rubrique{Entraînement}
\begin{blocLivre}{ex-13-6}\begin{exercice}[2]{Quatre angles autour d’un point}
Deux droites se coupent et un angle vaut $47^\circ$. Donner les trois autres mesures.
\label{fig-ex-13-6}\figureGeo{ang-6}
\end{exercice}\end{blocLivre}
\begin{blocLivre}{ex-13-7}\begin{exercice}[2]{Alternes internes}
Les droites $(d)$ et $(e)$ sont parallèles. L’angle marqué $a$ est alterne interne avec l’angle de $68^\circ$. Calculer $a$.
\label{fig-ex-13-7}\figureGeo{ang-7}
\end{exercice}\end{blocLivre}
\begin{blocLivre}{ex-13-8}\begin{exercice}[2]{Correspondants}
Les droites sont parallèles et deux angles correspondants sont marqués. L’un vaut $112^\circ$. Trouver l’autre.
\label{fig-ex-13-8}\figureGeo{ang-8}
\end{exercice}\end{blocLivre}
\begin{blocLivre}{ex-13-9}\begin{exercice}[2]{Prouver un parallélisme}
Les deux angles alternes internes indiqués mesurent chacun $72^\circ$. Que peut-on conclure sur les droites ?
\label{fig-ex-13-9}\figureGeo{ang-9}
\end{exercice}\end{blocLivre}
\begin{blocLivre}{ex-13-10}\begin{exercice}[2]{Exclure un parallélisme}
Deux angles correspondants mesurent $75^\circ$ et $78^\circ$. Les droites peuvent-elles être parallèles ?
\label{fig-ex-13-10}\figureGeo{ang-10}
\end{exercice}\end{blocLivre}

\rubrique{Problèmes}
\begin{blocLivre}{ex-13-11}\begin{exercice}[3]{Deux étapes}
Les droites sont parallèles. Au premier croisement, un angle obtus vaut $124^\circ$. Trouver l’angle aigu indiqué au second croisement.
\label{fig-ex-13-11}\figureGeo{ang-11}
\end{exercice}\end{blocLivre}
\begin{blocLivre}{ex-13-12}\begin{exercice}[3]{Équerres}
On utilise une équerre de $30^\circ$, $60^\circ$, $90^\circ$ et une de $45^\circ$, $45^\circ$, $90^\circ$. En juxtaposant un angle de $30^\circ$ et un de $45^\circ$, obtenir la mesure totale puis son supplément.
\label{fig-ex-13-12}\figureGeo{ang-12}
\end{exercice}\end{blocLivre}
\begin{blocLivre}{ex-13-13}\begin{exercice}[3]{Un angle à reconstituer}
Autour de O, trois angles consécutifs valent $110^\circ$, $95^\circ$ et $x$. Le dernier vaut $80^\circ$. Trouver $x$.
\label{fig-ex-13-13}\figureGeo{ang-13}
\end{exercice}\end{blocLivre}
\begin{blocLivre}{ex-13-14}\begin{exercice}[3]{Justifier sans mesurer}
Les droites $(d)$ et $(e)$ sont toutes deux perpendiculaires à la sécante, comme l’indique le codage. Une élève affirme qu’elles sont parallèles « parce qu’on le voit à la règle ». Proposer une justification correcte.
\label{fig-ex-13-14}\figureGeo{ang-14}
\end{exercice}\end{blocLivre}
\begin{blocLivre}{ex-13-15}\begin{exercice}[3]{Une bissectrice dans des parallèles}
Les droites sont parallèles. Un angle au premier croisement vaut $106^\circ$. Au second, on partage l’angle aigu correspondant en deux par une bissectrice. Trouver chaque moitié.
\label{fig-ex-13-15}\figureGeo{ang-15}
\end{exercice}\end{blocLivre}

\clearpage
